\slajdid{07-s003}
\begin{frame}[label=07-s003]
\gusttekst\setlength{\parskip}{5pt}
Digitalne, elektronske, komponente se mogu klasifikovati u dve velike grupe\par
\hspace{8mm}\textcolor{DEplava}{fiksne}\\\hspace{8mm}\textcolor{DEplava}{programabilne}\par
Pri ovome govorimo i mislimo o hardveru komponenti. Na primer: Računar je programabilan. Možemo da napišemo različite programe i da se računar izvođenjem programa ponaša na različite načine. Ali hardver računara je fiksan. Izvođenje programa ne menja hardverske veze u procesoru ili nekom drugom delu računara (zaboravite na trenutak memorije).\par\vspace{4mm}
Fiksne komponente – kada su proizvedene određena im je logička funkcija i ne može se menjati\par\vspace{4mm}
Programabilne komponente – korisnik može da ih prilagodi, programira, za određenu funkciju. Za sada to su neke logičke funkcije. Kombinacione mreže. Razvoj tehnologije je to omogućio. Da se sa kola malog i srednjeg stepena integracije koja su po pravilu fiksne komponente pređe na korišćenje komponenta višeg stepena integracije i koja mogu da se prilagode potrebi korisnika.
\end{frame}
